

\newcommand{\marked}[1]{#1}
\newcommand{\textmarked}[1]{#1}

%
%
\spis{Dwie tury, STV i~głosowanie strategiczne}
{Andrzej Kaczmarczyk,  Stanisław Szufa}
%
\wtyt{Dwie tury, STV i~głosowanie strategiczne}
\waut{Andrzej KACZMARCZYK* i~Stanisław SZUFA**}\marg{\setbox0=\hbox{**\kern1.2pt}\hangindent\wd0
\noindent\phantom{*}*\kern1.2ptDepartment of Computer Science,
The~University~of~Chicago 
\llap{**\kern1.2pt}Wydział~Informatyki,  Akademia~Górniczo-Hutnicza im.~Stanisława Staszica w~Krakowie\bigskip}
%
%

Wybory prezydenckie w~Polsce zwykle rozstrzygają się dopiero w~II turze.
%
Dochodzi do niej, gdy
w~I~turze żaden z~kandydatów nie zdobył ponad 50\%~ważnych głosów\marg{Choć ważność głosu to zagadnienie bardzo istotne, nie będziemy się na nim skupiać. Załóżmy, że wszystkie głosy, o~których mówimy, są ważne.\bigskip}.
Dla przypomnienia, II tura to kolejne głosowanie, w~którym o~zwycięstwo walczą
dwaj kandydaci z największym poparciem %, którzyzdobyli największe poparcie
w~I~turze.
Głosowanie dwuturowe jest względnie komfortowe dla wyborców.
Mogą oni przenieść swoje poparcie na innego kandydata w~II turze, na przykład gdy ich faworyt odpadł już w~I~turze.
Intuicyjnie więc rozsądny\marg{,,Rozsądny'' nie jest tu kontrastem dla ,,lekkomyślny''. Traktujemy słowo ,,rozsądny'' technicznie, jako maksymalizujący szansę zwycięstwa ulubionego kandydata.} wyborca powinien poprzeć w~I~turze swojego ulubionego kandydata.
Ale czy aby na pewno?
Zademonstrujemy,
że intuicja tutaj zawodzi. Są bowiem sytuacje, w~których lepiej poprzeć
nieco mniej lubianego kandydata i~zagwarantować mu zwycięstwo, niż zagłosować na swojego faworyta i~musieć pogodzić się ze
zwycięstwem bardzo nielubianego kandydata.
%


Rozważmy $63$ wyborców i~pięciu kandydatów \textmarked{$A$}, \textmarked{$B$}, \textmarked{$C$}, \textmarked{$D$} i~\textmarked{$E$}. W~przykładzie poniżej preferencje wyborców względem kandydatów podane są (czytając od lewej) od ulubionego kandydata. 
$$\jot0pt\begin{aligned}
  21\text{ wyborców}:& \quad %A \succ \lmarkelem{\marked{C}}{1} \succ \marked{E} \succ \rmarkelem{\marked{D}}{1} \succ B \highlight{L1-C}{R1-D}\\
  A \succ \tcbhighmath{C \succ E \succ D} \succ B\\
  19\text{ wyborców}:& \quad %B \succ \lmarkelem{\marked{C}}{2} \succ \marked{D} \succ \rmarkelem{\marked{E}}{2} \succ A \highlight{L2-C}{R2-E}\\
  B \succ \tcbhighmath{C \succ D \succ E} \succ A\\
  11\text{ wyborców}:& \quad %\lmarkelem{\marked{C}}{3} \succ \marked{E} \succ \rmarkelem{\marked{D}}{3} \succ B \succ A \highlight{L3-C}{R3-D}\\
  \tcbhighmath{C \succ E \succ D} \succ B \succ A\\
  7\text{ wyborców}:& \quad %\lmarkelem{\marked{D}}{4} \succ \marked{C} \succ \rmarkelem{\marked{E}}{4} \succ B \succ A \highlight{L4-D}{R4-E}\\
  \tcbhighmath{D \succ C \succ E} \succ B \succ A\\
  5\text{ wyborców}:& \quad %\lmarkelem{\marked{E}}{5} \succ \marked{C} \succ \rmarkelem{\marked{D}}{5} \succ A \succ B \highlight{L5-E}{R5-D}
        \tcbhighmath{E \succ C \succ D} \succ A \succ B 
\end{aligned}$$
Zauważmy, że gdyby wszyscy wyborcy oddali głos na ulubionego dostępnego w~danej turze kandydata,
to w~II turze $B$ pokonałby $A$ stosunkiem głosów $37$ do~$26$.

Przyjrzyjmy się jednak bliżej grupie $12$ wyborców przedstawionych jako ostatni w~powyższym zestawieniu. 
Po chwili refleksji zauważamy, że byliby oni dużo bardziej zadowoleni z~wyboru kandydata \textmarked{$C$}.
Rzeczywiście każdy taki wyborca uważa \textmarked{$C$} za swojego drugiego w~kolejności ulubionego kandydata.
Czy mogli oddać swoje głosy w~inny sposób, tak by zagwarantować kandydatowi~\textmarked{$C$} zwycięstwo? Zdecydowanie tak!
Zauważając, że ich ulubieni kandydaci \textmarked{$E$} oraz~\textmarked{$D$} są niezbyt popularni (popatrzmy, jak często wyprzedza
ich \textmarked{$C$}), wyborcy mogli oddać głos na~\textmarked{$C$}.\marg{Przykład udał nam się nawet lepiej niż się spodziewaliśmy! Poprzedni zwycięzca~$B$ nawet nie dostał się do~II~tury!}
Wtedy kandydat~\textmarked{$C$} z~$42$~głosami w~II~turze pokonałby kandydata $A$.

Skupmy się teraz na wyróżnionych kandydatach \textmarked{$C$}, \textmarked{$D$} oraz \textmarked{$E$}.
Tworzą oni spójny blok w~preferencjach każdego z~wyborców. Jak możemy to zinterpretować?
Kandydaci ci są na tyle podobni, że zawsze występują blisko siebie w~preferencjach wyborców, a~zarazem na tyle wyjątkowi,
że każdy inny kandydat jest albo bardziej preferowany od całej trójki, albo mniej preferowany od całej trójki.
Dodatkowo, skoro kandydaci zawsze występują blisko siebie w~preferencjach wyborców, intuicyjnie każdy wyborca będzie
,,podobnie'' zadowolony z~dowolnego z~tych kandydatów. Właśnie dlatego naukowcy, opisując takich kandydatów,
używają określenia \emph{klony}. W~praktyce (i~w~przybliżeniu) możemy myśleć o~klonach jako\marg{Klony mogą pochodzić na przykład z~opcji politycznych powstałych na skutek rozpadu jednej partii.} o~kandydatach o~bardzo zbliżonych poglądach lub
z~podobnych opcji politycznych.

Nasz przykład pokazał, że wybory dwuturowe czasem premiują
\emph{głosowanie strategiczne}, czyli takie, w~którym wyborcom 
bardziej opłaca się zagłosować niezgodnie z~własnymi preferencjami,
aby zwiększyć szansę na bardziej satysfakcjonujący wynik.
%
%
%
%
W~naszym przykładzie u~podstaw tego problemu leży podatność dwuturowej reguły wyborczej na \emph{klonowanie} kandydatów.
W~szczególnym przypadku możemy tę własność opisać tak: Wyobraźmy sobie, że klonujemy wygrywającego kandydata i~pozwalamy każdemu wyborcy doprecyzować, czy wolą kandydata oryginalnego, czy sklonowanego. Jeśli reguła jest podatna na klonowanie, to może się zdarzyć, że w~tak powstałych nowych wyborach nie wygra ani poprzednio wygrywający kandydat, ani jego klon.
%
%
%
%
%
%
%
%

Czy istnieje
racjonalna reguła, która nie jest podatna na klonowanie kandydatów?
Tak,\marg{W~szczególnym przypadku jednego zwycięzcy obok STV używane są też
nazwy {\em Instant-Runoff Voting}, {\em Ranked Choice Voting} czy {\em Alternative Vote}.} mechanika reguły STV (\emph{Single Transferrable Vote}, czyli pojedynczy głos przechodni) broni przed
zaobserwowaną sytuacją! W~przeciwieństwie do tradycyjnego głosowania dwuturowego,
które na każdym etapie wymaga od wyborców zaznaczenia tylko najbardziej
popieranego kandydata, reguła STV oczekuje uszeregowania wszystkich kandydatów.
Choć jest to trudniejsze od postawienia pojedynczego krzyżyka, to
%
%
\marg{Gdyby zastosować STV, zniknęłaby też potrzeba organizowania II tury wyborów.}%
taki sposób głosowania automatycznie umożliwia przechodzenie poparcia na kolejnego ulubionego kandydata w~przypadku,
gdy pierwszy nie może zwyciężyć ze względu na zbyt niskie poparcie. 

Jak więc działa mechanizm reguły STV przy wyborze jednego zwycięzcy?
%
%
%
%
%
%
Jest on podzielony na etapy, \marg{Liczba głosów potrzebna do zwycięstwa
nosi nazwę progu i~można go ustalić na kilka sposobów. My wybraliśmy ten najpopularniejszy (prawdopodobnie) i~zastosowaliśmy go do szczególnego przypadku wyboru jednego zwycięzcy.} a~w~każdym z~nich najpierw sprawdzamy, czy
jest taki kandydat, który zdobywa ponad połowę wszystkich głosów.
%
W~przypadku, gdy taki kandydat istnieje, to jest on zwycięzcą.
W~przeciwnym razie kandydat o~najniższym poparciu jest usuwany
z~grona kandydatów. Z~kolei głos każdego wyborcy, który popierał
usuniętego kandydata, przechodzi na kolejnego niewyeliminowanego
kandydata zgodnie z~listą preferencji wyborcy.
Następnie STV zaczyna następny etap i~znów szuka kandydata,
który zdobywa ponad połowę głosów,\marg{Uwaga! W~trakcie działania reguły mogą zdarzyć się remisy. W~takiej sytuacji zakładamy, że istnieje jakiś sposób ich rozstrzygania.}
%
kontynuując swoje działanie do momentu wyłonienia zwycięzcy.
Ostatecznie musi do tego dojść, gdyż w~najgorszym razie pozostanie
jedynie dwóch rozważanych kandydatów.

Przyjrzyjmy się zatem, jak zadziała STV na naszym przykładzie.
%
%
Na początku żaden z~kandydatów nie osiąga co najmniej $32$ głosów, więc reguła eliminuje najmniej popieranego
kandydata, \textmarked{$E$}. Głosy jego wyborców przechodzą na kandydata \textmarked{$C$}.
Dalej żaden z~kandydatów nie osiąga wymaganego poparcia, co powoduje usunięcie kandydata \textmarked{$D$} i~przejście jego głosów
również na kandydata \textmarked{$C$}. Teraz \textmarked{$C$} ma $23$ głosy. Do rozstrzygnięcia prowadzi wyeliminowanie następnego
najsłabszego kandydata, $B$, którego wyborcy przenoszą swoje poparcie na kandydata \textmarked{$C$}, a~ten z~$42$ głosami  zostaje zwycięzcą.

Jak pokazaliśmy, reguła STV doprowadziła kandydata \marked{$C$} do zwycięstwa, pozwalając wyborcom
głosować całkowicie szczerze.\marg{Reguła STV ma też inne atuty, takie jak łatwość dostosowania jej do wyborów z~wieloma zwycięzcami w~sposób zapewniający proporcjonalną reprezentację wyborców.} Intuicyjnie jest to zasługa faktu, że sam mechanizm reguły naturalnie przenosi poparcie wyborców,
których ulubiony kandydat jest eliminowany.
%
Wyborca efektywnie wyraża
swój głos w~formie ,,jeśli nie kandydat $X$, to $Y$, a~jeśli nie $Y$, to $Z$'', zamiast wskazania tylko $X$.
Wiąże się to jednak z~kosztem wynikającym z~potrzeby uszeregowania kandydatów.
Choć w~małych wyborach uszeregowanie kilku kandydatów nie wydaje się trudnym zadaniem, to kiedy mowa
o~kilkudziesięciu kandydatach, łatwo o~pomyłkę, która może unieważnić głos.
Dodatkowo samo uświadomienie sobie własnych preferencji może być bardzo trudne.

Należy zaznaczyć, że reguła STV z~powodzeniem stosowana jest w~praktyce. Przykładowo zgodnie z~nią wybierany jest prezydent Irlandii,
reprezentanci izby niższej australijskiego parlamentu czy przedstawiciele samorządów lokalnych w~Szkocji. W~przywołanym przykładzie Australii STV wybiera po jednym kandydacie w~każdym okręgu. W~Szkocji, w~zależności od okręgu, wybieranych jest od jednego do pięciu kandydatów.

Uczulamy jednak na wyciąganie zbyt pochopnych wniosków.
Reguła STV\marg{Wszystkie rozsądne reguły wyborcze mają swoje ciemne strony. Mówią o~tym twierdzenia Arrowa i~Gibbarda--Satterthwaite'a, z~którymi zaznajomienie się polecamy Czytelnikom Zainteresowanym.}  ma również swoje ciemne strony.
Jedna z~nich jest szczególnie kontrowersyjna.
Polega ona na tym, że może się zdarzyć, że wygrywający kandydat~$X$ przegra wybory w~sytuacji, gdy wyborcy postanowią przesunąć go wyżej na swojej liście preferencji.
To znaczy, że wyborca chcący bardziej wesprzeć $X$ działa na jego niekorzyść. Spójrzmy na przykład poniżej, w~którym reguła STV wybiera kandydata $X$, po drodze eliminując $Y$.
\begin{align*}
10\text{ wyborców}:& \quad X \succ Y \succ Z\\
9\text{ wyborców}:& \quad Y \succ X \succ Z\\
8\text{ wyborców}:& \quad Z \succ Y \succ X\\
2\text{ wyborców}:& \quad Z \succ X \succ Y
\end{align*}
Gdyby jednak dwóch ostatnich wyborców postanowiło, że
przesuną $X$ wyżej na swoich listach preferencji,
to kandydat $X$ przegrałby na rzecz $Y$!
To~ważna obserwacja, ponieważ prowadzi ona czasem
do sporów sądowych o~konstytucyjność reguł wyborczych
podatnych na opisaną anomalię. W~Niemczech odpowiedni
sąd uznał ją za przesłankę wystarczającą 
do uznania rozpatrywanej przez niego reguły za niekonstytucyjną
i~zakazał jej stosowania.
Przeciwnicy tego punktu widzenia argumentują, że w~praktyce ten problem jest pomijalny: wydaje się występować rzadko oraz nie wpływa ani na zachowania wyborców, ani na zachowania kandydatów -- czego nie można powiedzieć o~problemie klonów. 

Czytelnikom Zainteresowanym dalszym pogłębieniem tematyki głosowania 
strategicznego polecamy poniższe publikacje:
\begin{enumerate}
    \item Doron, G. i~Kronick, R.: ``Single Transferrable Vote: An Example of a~Perverse Social Choice Function'', \emph{American Journal of Political Science}. 1977, T. 31, nr 2, s. 303--311.
    \item Taylor, A. D. ``{Social Choice and the Mathematics of Manipulation}'',  Cambridge University Press, Cambridge 2005.
\end{enumerate}

\vfill
\input 06-kat-o

\eject